\documentclass[11pt,letterpaper]{article}

\usepackage[margin=2.5cm]{geometry}
\usepackage{amsmath}
\usepackage{amssymb}
\usepackage{bm}
\usepackage{booktabs}
\usepackage{array}
\usepackage{ragged2e}
\usepackage{seqsplit}
\usepackage{enumitem}
\usepackage{xcolor}
\usepackage[hidelinks]{hyperref}
\usepackage{caption}
\usepackage{float}
\usepackage{graphicx}
\usepackage{setspace}
\setstretch{1.15}

\DeclareGraphicsExtensions{.pdf}
\graphicspath{{figs/}}
\setlength{\emergencystretch}{3em}
\newcolumntype{L}[1]{>{\RaggedRight\arraybackslash}p{#1}}

%----------------------------------------------------------------------
%  Text helpers
%----------------------------------------------------------------------
\newcommand{\file}[1]{\texttt{\seqsplit{#1}}}
\newcommand{\code}[1]{\texttt{#1}}
\newcommand{\SCALAR}{\textbf{scalar}}
\newcommand{\ToClaude}[1]{[\textbf{CLAUDE}: {#1}]}

%----------------------------------------------------------------------
%  Notation: the same as tex/analysis_chain.tex.  Images bold, tables
%  blackboard bold, scalars plain.  \overline = mean over a region,
%  \widetilde = smoothed, \widehat = fitted or estimated.
%----------------------------------------------------------------------
\newcommand{\mR}{\mathbf{R}}   % post-image reflectance vs gray chips
\newcommand{\mS}{\mathbf{S}}   % pre-image reflectance, lined up
\newcommand{\mA}{\mathbf{A}}   % post-image absorbance
\newcommand{\mB}{\mathbf{B}}   % pre-image absorbance, lined up
\newcommand{\mC}{\mathbf{C}}   % stain absorbance
\newcommand{\mD}{\mathbf{D}}   % error image (this study only)
\newcommand{\mF}{\mathbf{F}}   % dark-corrected signal
\newcommand{\mG}{\mathbf{G}}   % gray-chip level map
\newcommand{\mP}{\mathbf{P}}   % flat-field polynomial shape
\newcommand{\mM}{\mathbf{M}}   % region indicator (1 inside, 0 outside)
\newcommand{\mX}{\mathbf{X}}   % any image
\newcommand{\tB}{\widetilde{\mathbf{B}}}
\newcommand{\hB}{\widehat{\mathbf{B}}}
\newcommand{\tD}{\widetilde{\mathbf{D}}}
\newcommand{\tX}{\widetilde{\mathbf{X}}}
\newcommand{\Sown}{\mathbf{S}_{\text{own}}}
\newcommand{\Bown}{\mathbf{B}_{\text{own}}}
\newcommand{\Rwhite}{R_{\text{white}}}
\newcommand{\bw}{b_{\text{w}}}
\newcommand{\Cbar}{\overline{C}}
\newcommand{\reg}[2]{\underset{\text{#1}}{#2}}
\newcommand{\avg}[2]{\overline{\reg{#1}{#2}}}
\newcommand{\MA}{\mathcal{A}}
\newcommand{\MB}{\mathcal{B}}
\newcommand{\MC}{\mathcal{C}}
\newcommand{\MD}{\mathcal{D}}

\title{Shading removal with the pre-image\\
\large Four methods, how lining up changes the pre-image,\\
and what the single white value $\Rwhite$ does}
\author{Colorimeter project --- pad analysis chain}
\date{2026-10-03 (rewritten in the notation of \file{analysis\_chain.tex};
first version 2026-10-02).  Branch \file{preimage-capture}, commit
\file{4962604} plus uncommitted changes to
\file{desktop/padanalysis/shading\_review.py}.}

\begin{document}
\maketitle

\begin{abstract}
\noindent
A photo of the indicator pad is darker in some places than others even
where there is no stain.  The pad's edge casts a shadow, and the paper's
fibers make a fine pixel-to-pixel pattern.  The pre-image (the same strip,
photographed before exposure to arsine) holds the same shadow and much of
the same pattern, so taking it out should leave only the stain.  This note
compares four ways of doing that on one pad (C1032, royal\_blue.B).  It
writes each as one equation, $\mC = \mA - \hB$, and shows where a single
number is applied to every pixel.  Three findings:
\begin{itemize}[nosep]
  \item No method wins everywhere: smoothing the pre-image is best inside
    the pad, using it pixel by pixel is best at the rim.
  \item The single white value $\Rwhite$ cancels in three methods and is
    right only more than 2 pixels in from the pad's edge.
  \item Lining up the pre-image on the post-image averages 4 pixels, and
    that lowers its pixel-to-pixel scatter to 0.61 of its true size
    (median over 18 pads and 6 channels; range 0.43--0.96, set by where
    each pad happened to land).  So ``pixel by pixel'' is really a
    smoothing whose strength changes from pad to pad.
\end{itemize}
The chain is not changed by this note.
\end{abstract}

\tableofcontents
\newpage

%======================================================================
\section{The question, the decision, and when to stop}

\begin{description}[leftmargin=2em]
  \item[Question.] Which way of taking out the pre-image leaves the
    corrected pad flattest where there is no stain, without changing the
    stain?
  \item[Decision it could change.] Step 6 of the chain
    (\file{analysis\_chain.tex}): how $\hB$ is made.
  \item[When to stop.] When the methods have been compared on one pad
    with figures and numbers (this note), then on all 18 pads and six
    channels as a table.  Choose a method only after the second step.
\end{description}

%======================================================================
\section{Notation}
\label{sec:notation}

The notation is that of \file{analysis\_chain.tex}, Section 2.  In short:
\textbf{bold uppercase} is an image (one value per pixel, one channel);
plain is a single number; $A_{ij}$ is the pixel in row $i$, column $j$; a
region's name is written under a symbol; an overline is a mean over a
region, a tilde means smoothed, a hat means estimated or fitted.  Arithmetic
between images is pixel by pixel.  Scalars that are applied to many pixels
are flagged in bold, because each assumes those pixels are alike.

\medskip
\begin{tabular}{@{}L{2.3cm}L{1.6cm}L{10.7cm}@{}}
\toprule
Symbol & Kind & Meaning \\
\midrule
$\mR$ & image & post-image reflectance: dark-corrected signal divided by
  the gray-chip level at the same pixel.  Bigger = lighter.  The pad reads
  about 2 (it is about twice as bright as the gray chips). \\
$\mS$ & image & the pre-image, the same steps, lined up on $\mR$'s grid. \\
$\Sown$ & image & the pre-image before lining up, on its own grid. \\
$\mA$, $\mB$, $\Bown$ & image & absorbance: $-\log_{10}$ of $\mR$, $\mS$,
  $\Sown$.  Bigger = darker.  0.01 in absorbance is about 2.3\,\% in
  reflectance. \\
$\tX$ & image & any image $\mX$ smoothed over width $\sigma$, pad pixels
  only (Section~\ref{sec:smooth}). \\
$\hB$ & image & the estimated shading: what a method takes out. \\
$\mC$ & image & the stain: $\mC = \mA - \hB$. \\
$\Rwhite$ & \SCALAR & the single white value of method $\MA$: the median of
  $\Sown$ more than 6 pixels in from the pad's edge.  One per channel per
  strip.  Not part of the chain (\file{analysis\_chain.tex}). \\
$\bw$ & \SCALAR & the same in absorbance, $\bw = -\log_{10}\Rwhite$. \\
$\lambda$ & \SCALAR & method $\MD$: the weight on $\mB$; $1 - \lambda$ goes
  to $\tB$.  $\hat\lambda$ is its fitted value. \\
$\Cbar$ & \SCALAR & one per method: the median of $\mC$ over the dye-free
  inner pad.  Used only to measure error. \\
$\mD$, $\tD$ & image & the error, $\mD = \mC - \Cbar$, and its smooth
  part (Section~\ref{sec:error}).  This study only. \\
\bottomrule
\end{tabular}

\medskip
\noindent Regions of the post-image pad:

\medskip
\begin{tabular}{@{}L{3.2cm}L{11.4cm}@{}}
\toprule
pad & the fitted pad rectangle, hot pixels left out: 2215 pixels on C1032. \\
rim, inner & pad pixels within, and more than, $w = 4$ pixels of the pad's
  edge. \\
dot & pad pixels inside the dot's fitted edge (dyed): 280 pixels. \\
free & pad pixels at least 6 pixels outside the dot's edge (dye-free).
  Split into the dye-free inner pad (760 pixels) and the dye-free rim (686
  pixels). \\
g1, g2 & the left and right gray chips. \\
\bottomrule
\end{tabular}

\medskip
\noindent The dot's size and shape on these samples are not final (the
exposure geometry will change), so nothing here depends on the dot being
round.

%======================================================================
\section{Where \texorpdfstring{$\mR$ and $\mS$}{R and S} come from}
\label{sec:chain}

Steps 1--4 of \file{analysis\_chain.tex}, in brief.  Each photo:
\begin{enumerate}
  \item \textbf{Dark level.}  $F_{ij} = E_{ij} - e$: the lit layer minus
    $e$, a \SCALAR{} per Bayer layer per photo (the mean of one dark
    frame).
  \item \textbf{Gray-chip level.}  $R_{ij} = F_{ij} / G_{ij}$.  $\mG$ is
    the flat-field polynomial $\mP$ (a fixed shape, measured once) scaled to
    match the two gray chips of this same photo.  The \SCALAR{} mean signals
    $\avg{g1}{\mF}$ and $\avg{g2}{\mF}$ are taken on each chip.  The chips
    lie on opposite sides of the pad, and the scale is blended in a straight
    line between the chips' centroids, so that $\mG$ matches each chip there.
    That blend is in the column direction only; nothing is measured or
    adjusted in the row direction.
\end{enumerate}
The pre-image goes through both steps in its own frame, with its own gray
chips (the light belongs to the stage, not the strip), giving $\Sown$.
Then it is lined up on the post-image's grid, giving $\mS$: the two fitted
pad rectangles give a shift and scale (\code{analysis\_1.pre\_to\_post}),
and each lined-up value is a weighted average of the 4 nearest pre-image
pixels (\code{\_register}).  The next section measures what that averaging
does.

%======================================================================
\section{Lining up lowers the pre-image's pixel-to-pixel scatter}
\label{sec:lineup}

\subsection{Why it matters}

A lined-up value is a weighted average of 4 pre-image pixels, so it varies
less from pixel to pixel than the photo does.  That changes any noise or
texture measured on $\mS$, and it makes method $\MC$ (``pixel by pixel'')
partly a smoothing.  How much depends on where the post-image's pixels
fall between the pre-image's: not at all if they land on whole pixels, most
if they land half-way.

\subsection{How it is measured}

Three versions of the pre-image's absorbance are compared on the inner pad
(so the edge shadow does not count):
\begin{itemize}[nosep]
  \item $\Bown$, on the pre-image's own grid;
  \item lined up by taking the single \emph{nearest} pre-image pixel (no
    averaging);
  \item lined up as the chain does (\emph{bilinear}: the 4-pixel average).
\end{itemize}
For each, the difference between every pixel and its right-hand neighbor
(and, separately, the one below) is taken:
\begin{equation}
  d_{ij} = \frac{X_{i,j+1} - X_{ij}}{\sqrt 2}.
\end{equation}
Differences cancel anything smooth, so what is left changes from one pixel
to the next.  On the pre-image's pad that is mostly the paper's fiber
texture, with camera noise (about 0.002) a small part.  Dividing by
$\sqrt 2$ makes the spread of $d$ the spread of one pixel, if pixels are
independent.  The spread is the robust one, $1.4826 \times$ the median
distance from the median, so a speck of dust does not decide it
(\code{duotone.\_pair\_noise}).

\noindent\textbf{What pure camera noise would do.}  Let $f_x$ and $f_y$ be
how far a lined-up pixel falls between pre-image columns and rows (0 to 1;
scalars for each pixel).  For noise that is independent from pixel to
pixel, the 4-pixel average scales the spread of a difference along a row by
\begin{equation}
  \sqrt{\tfrac12\left[(1-f_x)^2 + (1-2f_x)^2 + f_x^2\right]\,
        \left[(1-f_y)^2 + f_y^2\right]},
\end{equation}
and along a column the same with $f_x$ and $f_y$ swapped.  At $f = 0$ the
factor is 1; at $f_x = 0.5$ the first bracket gives 0.5.  This is the
``predicted'' ratio below (its root-mean-square over the inner pad).

\subsection{Results}

\textbf{C1032 royal\_blue.B} (fractional position: rows 0.83, columns 0.05,
so mostly a row shift):

\medskip
\begin{center}
\begin{tabular}{@{}lrrrrr@{}}
\toprule
differences & own grid & nearest & bilinear & ratio bilinear/own & predicted \\
\midrule
along rows    & 0.0067 & 0.0067 & 0.0054 & 0.80 & 0.81 \\
along columns & 0.0068 & 0.0068 & 0.0049 & 0.72 & 0.74 \\
\bottomrule
\end{tabular}
\end{center}

\noindent\textbf{All 18 pads, six channels} (108 pairs; median, with range):
own grid 0.0067, nearest 0.0067, bilinear 0.0041.  Ratio bilinear/own
0.61 along rows and 0.61 along columns, range 0.43--0.96.  Predicted
0.58 and 0.57.  Measured is above predicted by a median of 0.04, because
the fiber texture spans more than one pixel and is averaged less than
independent noise would be.  C1032 is
among the milder cases (C1023 and C1033 are milder); C1017, C1019, C1020,
C1022, C1025, C1030, C1034 and C1035 fall near half a pixel and lose about
half their scatter.

\begin{figure}[h!]
    \centering
    \includegraphics[width=\textwidth]{shading-R4a}
    \caption{
    royal\_blue.B pre-image, inner pad: $\mB$ on its own grid, lined up by
    nearest pixel, and lined up by the 4-pixel (bilinear) average, each
    minus its median ($\pm 0.02$, bigger = darker); right, histograms of
    the neighbor differences along rows.  Top, C1032 (lands nearly on whole
    columns: mild averaging); bottom, C1020 (lands half-way: the scatter
    halves).  Made by \code{shading\_review.show\_lineup\_noise} (step R4a).
    }
\label{fig:shading-R4a}
\end{figure}

\begin{figure}[h!]
    \centering
    \includegraphics[width=0.62\textwidth]{shading-R4b}
    \caption{
    All 18 pads and six channels: the measured lowering of the
    pixel-to-pixel scatter (robust spread lined up by bilinear average,
    divided by the spread on the own grid) against the lowering predicted
    for independent pixel noise from each pad's fractional position.  1 =
    not lowered.  Made by \code{shading\_review.lineup\_noise\_table} and
    \code{show\_lineup\_ratio} (step R4b).
    }
\label{fig:shading-R4b}
\end{figure}

\subsection{What follows from it}

\begin{enumerate}
  \item \textbf{Measure noise and texture on the own grid}, never on the
    lined-up pre-image.  The yellow/brown code already does
    (\code{duotone.PadCube} reads the pre-image with \code{register=False}).
  \item \textbf{Method $\MC$ is not purely pixel by pixel.}  Its $\mB$ is a
    4-pixel average whose strength is set by where the pad landed, from
    almost none (C1033) to a half-pixel average (C1017, C1020).  So the
    methods' ranking, and the best $\lambda$ in method $\MD$, can differ from
    pad to pad for a reason unrelated to the paper.  This fits Bill's note
    that $\lambda$ may have to be set for each photo from measures of its
    noise and of how the photos line up.
  \item \textbf{Nearest-pixel lining up keeps the scatter} but places
    values up to half a pixel wrong, which matters at the edge shadow.
  \item \textbf{The C1032 results below are a mild case} (ratio 0.80 along
    rows).  Section~\ref{sec:open} lists what this suggests for the next
    step.
\end{enumerate}

%======================================================================
\section{The single white value \texorpdfstring{$\Rwhite$}{R\_white}
(method \texorpdfstring{$\MA$}{A})}
\label{sec:white}

\subsection{What it is}

\begin{align}
  \Rwhite &= \operatorname{median}\bigl\{\,(S_\text{own})_{ij} :
    \text{more than 6 px in from the pre-image pad's edge} \,\bigr\}
    && \text{(\SCALAR)} \\
  \bw &= -\log_{10}\Rwhite && \text{(\SCALAR)}
\end{align}
Method $\MA$ divides every post-image pixel by it; in absorbance it takes
out the same $\bw$ at every pixel:
\begin{equation}
  C_{ij} = A_{ij} - \bw = -\log_{10}\frac{R_{ij}}{\Rwhite}.
\end{equation}

\noindent\textbf{Worked example} (C1032 royal\_blue.B).  $\Rwhite = 2.1888$,
$\bw = -0.3402$.  The dye-free pixel at row 55, column 78 has
$R = 2.1327$, $A = -0.3289$, so $C = -0.3289 + 0.3402 = +0.0113$: 2.6\,\%
darker than the white value.

\subsection{If it cancels, what use is it?}

In methods $\MB$, $\MC$ and $\MD$ it cancels exactly
(Section~\ref{sec:cancel}): there the pre-image gives every pixel its own
white, $10^{-\hat B_{ij}}$.  In method $\MA$, the code's current behavior,
it does not cancel, not even in the yellow/brown split, which removes only
the average over channels at each pixel (\code{duotone.contrast}).  Its job
in $\MA$ is to make $\mC = 0$ mean ``as the unreacted pad'' in every
channel; without it, the paper's own color (how it differs from the gray
chips from one LED to the next) would be read as stain.  The size of that,
on C1032:

\medskip
\begin{center}
\begin{tabular}{@{}lcc@{}}
\toprule
channel & $\Rwhite$ & $-\bw = \log_{10}\Rwhite$ \\
\midrule
violet.B      & 1.8219 & 0.2605 \\
royal\_blue.B & 2.1888 & 0.3402 \\
cyan.B        & 2.0233 & 0.3061 \\
green.Gb      & 1.9398 & 0.2878 \\
green.Gr      & 1.9382 & 0.2874 \\
photo\_red.R  & 1.6101 & 0.2069 \\
\bottomrule
\end{tabular}
\end{center}

\noindent These differ by up to 0.13 between channels, about as much as the
palest stained pixels.  Mean or median makes no difference: over the region
it is taken from, the mean of $\Sown$ is 1.0002 times the median.  Over the
whole pad, rim included, the mean would be 2.4\,\% lower.

\subsection{Is it the right value everywhere on the pad?}

No.  Figure~\ref{fig:shading-R3d} shows each pre-image pad pixel against
$\Rwhite$, as $\Bown - \bw = -\log_{10}(\Sown/\Rwhite)$, by depth into the
pad:

\medskip
\begin{center}
\begin{tabular}{@{}lrrrrr@{}}
\toprule
depth into the pad & pixels & median & sd & 95th pct & median, in reflectance \\
\midrule
over 6 px ($\Rwhite$ taken here) & 1226 & 0.0000 & 0.0072 & +0.0119 & 0 \\
2 to 6 px    & 625 & +0.0035 & 0.0095 & +0.0207 & 0.8\,\% darker \\
1 to 2 px    & 176 & +0.0377 & 0.0191 & +0.0632 & 8.3\,\% darker \\
outermost pixel & 184 & +0.0686 & 0.0496 & +0.1886 & 14.6\,\% darker \\
\bottomrule
\end{tabular}
\end{center}

\noindent (The darkest outermost pixel is $+0.239$, 42\,\% darker.)  More
than 2 pixels in, $\Rwhite$ is right on average (within 0.0035) but each
pixel differs from it by the paper's texture (sd 0.0072, about 1.7\,\%).
In the outer 2 pixels it is wrong by 8--15\,\% typically: method $\MA$ reads
the edge shadow as stain.  The dot, on these samples, lies inside the
region $\Rwhite$ is taken from.

\begin{figure}[h!]
    \centering
    \includegraphics[width=\textwidth]{shading-R3d}
    \caption{
    C1032 royal\_blue.B pre-image, own grid: how far each pad pixel is from
    the single white value, as $\Bown - \bw$ (0 = $\Rwhite$; bigger =
    darker).  Left: the map; the black line encloses the region $\Rwhite$
    is the median of.  Right: histograms by depth into the pad, sharing the
    $x$ axis; black line = $\Rwhite$, red dashed = that band's median.
    Made by \code{shading\_review.show\_white} (step R3d).
    }
\label{fig:shading-R3d}
\end{figure}

%======================================================================
\section{The four methods, as equations}
\label{sec:methods}

\subsection{One equation, four choices of \texorpdfstring{$\hB$}{B-hat}}

Every method computes the stain the same way and differs only in the
estimated shading $\hB$ it takes out:
\begin{equation}
  C_{ij} = A_{ij} - \hat B_{ij}.
\end{equation}

\medskip
\begin{tabular}{@{}L{3.4cm}L{7.4cm}L{3.6cm}@{}}
\toprule
method & $\hat B_{ij}$ & scalars \\
\midrule
$\MA$ (code now) & $\bw$, the same at every pixel & $\Rwhite$ \\
$\MB$1, $\MB$2 & $\tilde B_{ij}$, $\mB$ smoothed over $\sigma = 1$ or 2 px & $\sigma$ \\
$\MC$ & $B_{ij}$, the lined-up pre-image itself & --- \\
$\MD$ & $\lambda B_{ij} + (1-\lambda)\tilde B_{ij}$, $\sigma = 1$ px & $\lambda$, $\sigma$ \\
\bottomrule
\end{tabular}

\medskip
\noindent In reflectance each is a division by a per-pixel white,
$10^{-\hat B_{ij}}$; for $\MA$ that white is the single number $\Rwhite$.

\subsection{Why a constant cancels in \texorpdfstring{$\MB$, $\MC$ and $\MD$}{B, C and D}}
\label{sec:cancel}

Add the same number $c$ to every pixel of both $\mA$ and $\mB$ (for
example, by measuring both against $\Rwhite$ rather than against the gray
chips, $c = -\bw$).  Then:
\begin{itemize}
  \item $\MC$: $\hat B$ gains $c$, so $C = (A + c) - (B + c) = A - B$.
    Unchanged.
  \item $\MB$: smoothing is a weighted average whose weights add up to 1
    (Section~\ref{sec:smooth}), so $\tilde B$ gains exactly $c$ too.
    Unchanged.
  \item $\MD$: $\hat B$ gains $\lambda c + (1-\lambda) c = c$.  Unchanged.
    \textbf{The weights summing to 1 is what makes this work.}  Fitted
    freely, the weights on $\tB$ and on $\mB - \tB$ come out at 0.94 and
    0.87; then $1 - 0.94 = 0.06$ of the pad's own level against the gray
    chips ($\approx 0.34$) would stay in $\mC$: 0.020 at every pixel.
  \item $\MA$: $\hat B = \bw$ is a fixed number, so $C$ gains $c$.  Here the
    reference matters.
\end{itemize}

\subsection{The smoothing, exactly}
\label{sec:smooth}

\begin{equation}
  \tilde X_{ij} = \frac{\sum_{i',j'} \kappa_\sigma(i - i', j - j')\,
                        M_{i'j'}\, X_{i'j'}}
                       {\sum_{i',j'} \kappa_\sigma(i - i', j - j')\, M_{i'j'}}
\end{equation}
where $\kappa_\sigma$ is a bell-shaped weight of width $\sigma$ pixels (a
Gaussian with standard deviation $\sigma$), the sums run over pixels near
$(i, j)$, and $M_{i'j'} = 1$ on the pad, 0 elsewhere.  Dividing by the sum
of the weights makes them add up to 1 at every pixel.  Near the pad's edge
only pad pixels count, so the paper outside does not leak in; but a pixel at
the edge is then averaged only with pixels further in, which are lighter.
\textbf{So smoothing pulls the edge shadow inward} and makes the outermost
pixel look lighter than it is (Section~\ref{sec:results}).

\subsection{How \texorpdfstring{$\lambda$}{lambda} is fitted}

Method $\MD$ is $\mC = \mA - \tB - \lambda(\mB - \tB)$.  $\hat\lambda$ is the
value that makes the sum of squared leftovers smallest over the dye-free
zone (an overall offset is allowed for in the fit, then thrown away):
\begin{equation}
  \hat\lambda = \frac{\sum (f - \bar f)(y - \bar y)}{\sum (f - \bar f)^2},
  \qquad f = B_{ij} - \tilde B_{ij},\quad y = A_{ij} - \tilde B_{ij},
\end{equation}
summed over the dye-free pixels, with $\bar f$, $\bar y$ their means there.
On C1032: $\hat\lambda = 0.81$, one \SCALAR{} for the whole pad.
\textbf{Fitted on each part alone, it is 0.24 on the dye-free inner pad and
0.85 on the dye-free rim.}  The rim's error is about six times larger, so
it sets the pooled value.

\subsection{Worked example: two pixels, every method}

Both pixels are dye-free.  The inner one is 8.8 px in from the pad's edge
(row 55, column 78); the rim one is 0.8 px in (row 66, column 70).
$\Rwhite = 2.1888$, $\bw = -0.3402$.

\medskip
\begin{center}
\begin{tabular}{@{}lrr@{}}
\toprule
 & inner pixel & rim pixel \\
\midrule
$R$ & 2.1327 & 1.5559 \\
$S$ & 2.1939 & 1.5702 \\
$A = -\log_{10} R$ & $-0.3289$ & $-0.1920$ \\
$B = -\log_{10} S$ & $-0.3412$ & $-0.1960$ \\
$\tilde B$, $\sigma = 1$ & $-0.3418$ & $-0.2431$ \\
\midrule
\multicolumn{3}{@{}l}{$C = A - \hat B$ for each method:} \\
$\MA$: $\hat B = \bw$                            & +0.0113 & +0.1482 \\
$\MB$1: $\hat B = \tilde B$, $\sigma = 1$        & +0.0129 & +0.0511 \\
$\MB$2: $\hat B = \tilde B$, $\sigma = 2$        & +0.0117 & +0.0853 \\
$\MC$: $\hat B = B$                              & +0.0123 & +0.0040 \\
$\MD$: $\hat B = 0.81\,B + 0.19\,\tilde B$       & +0.0124 & +0.0129 \\
\bottomrule
\end{tabular}
\end{center}

\noindent Check of $\MC$ on the rim pixel, from reflectance:
$R/S = 1.5559/1.5702 = 0.9909$, $-\log_{10}(0.9909) = +0.0040$; no
$\Rwhite$ needed.  Check of $\MD$ on the rim pixel:
$\hat B = 0.81(-0.1960) + 0.19(-0.2431) = -0.2049$, and
$C = -0.1920 + 0.2049 = +0.0129$.

\noindent On the inner pixel all methods agree within 0.0016, under the
paper's texture (0.007).  On the rim pixel they differ by 0.14: the pad is
0.14 darker there than inside in both photos; $\MA$ leaves all of it, $\MB$1
leaves a third (smoothing has lightened $\tilde B$ there from $-0.196$ to
$-0.243$), and $\MC$ removes it.

%======================================================================
\section{How the error is measured}
\label{sec:error}

The truth is not known at any pixel, so the error has to be defined.  The
rule: \emph{the corrected dye-free pad should be flat, at the level of the
dye-free inner pad.}  For each method:
\begin{align}
  \Cbar &= \operatorname{median}\{C_{ij} : \text{dye-free inner pad}\}
    && \text{(\SCALAR, one per method)} \\
  D_{ij} &= C_{ij} - \Cbar && \text{(the error; image, whole pad)}
\end{align}
The error is split into a smooth part $\tD$ ($\mD$ smoothed with
$\sigma = 2$ px) and a fine part $\mD - \tD$.
\begin{itemize}
  \item $\Cbar$ only sets where 0 is.  It is subtracted from every pad
    pixel, so $\Rwhite$ has no effect on $\mD$ in any method.  But because it
    centers each method's inner pad on 0, a method that shifts the whole
    inner pad does not show it in $\mD$: $\MB$2 does ($\Cbar = 0.0027$
    against 0.0044--0.0048 for the others).
  \item $\tD$ is shadow left over (positive) or overdone (negative).  2 px
    is wider than the fiber texture (1--2 px) and narrower than the edge
    shadow (2--3 px).
  \item $\mD - \tD$ is everything that changes within about 2 px: camera
    noise, fiber texture not cancelled, one-pixel mismatches at the pad's
    edge, and, around the dot, the stain's own steep edge.
  \item Inside the dot, $\mD$ is stain, not error.  Error numbers are for
    the dye-free zone only; for the dot only the mean of $\mC$ is compared.
  \item $\tD$ is smoothed over the whole pad, so a little of the stain's
    fade leaks into it near the start of the dye-free zone: on average
    +0.0004 to +0.001 within 3 px of that line, at most +0.004.
\end{itemize}

%======================================================================
\section{Results on C1032 royal\_blue.B}
\label{sec:results}

\subsection{Orientation}

\begin{figure}[h!]
    \centering
    \includegraphics[width=\textwidth]{shading-R2a}
    \caption{
    C1032 royal\_blue.B: $\mC$ by method $\MA$ ($\mA - \bw$), the lined-up
    pre-image ($\mB - \bw$), and $\mC$ by method $\MC$ ($\mA - \mB$).  Top
    row: a scale that shows the dot.  Bottom row: a tight scale ($\pm 0.04$)
    that shows shading.  Black line = dot edge; gray line = start of the
    dye-free zone.  Bigger = darker.  Made by
    \code{shading\_review.show\_r2\_maps} (step R2a).
    }
\label{fig:shading-R2a}
\end{figure}

\subsection{One row, every method}

Figure~\ref{fig:shading-R3e} follows one dye-free row (row 82) and one row
through the dot (row 66) across the pad:
\begin{itemize}
  \item Panel 1: the edge shadow is in both photos (left end: $\mA$ rises
    0.12, $\mB$ 0.14 above the inner pad).  Smoothing lowers the outermost
    value and raises its neighbors: the shadow pulled inward.
  \item Panel 2: $\hB$ for each method; $\MA$'s is the flat line $\bw$.
  \item Panel 3, left end: $\MA$ leaves the shadow (+0.12); $\MB$1 and
    $\MB$2 leave part of it in the outermost pixel and overdo the next one
    in ($-0.01$); $\MC$ and $\MD$ come closest.
  \item Panel 3, right end: here the post-image is darker at the edge than
    the pre-image, so every method leaves +0.04 to +0.06.  The mismatch
    between the two photos' edges goes both ways.
  \item Panel 4: through the dot, all methods lie on top of each other;
    they differ only at the pad's edges.
\end{itemize}

\begin{figure}[h!]
    \centering
    \includegraphics[width=0.78\textwidth]{shading-R3e}
    \caption{
    C1032 royal\_blue.B, one row across the pad; panels share the column
    axis; the rim (within 4 px of the pad's edge) is shaded gray.  1:
    $\mA$, $\mB$ and $\tB$ against the gray chips; the gray line is $\bw$.
    2: $\hB$, what each method takes out.  3 and 4: $\mC = \mA - \hB$ on the
    dye-free row and on the row through the dot.  Bigger = darker.  Made by
    \code{shading\_review.show\_r3\_profile} (step R3e).
    }
\label{fig:shading-R3e}
\end{figure}

\subsection{Maps of the error}

\begin{figure}[h!]
    \centering
    \includegraphics[width=\textwidth]{shading-R3a}
    \caption{
    C1032 royal\_blue.B: where each method leaves its error.  One column per
    method; the scalars each uses are printed under its name.  Rows: 1,
    $\hB - \bw$ (what is taken out, relative to method $\MA$'s single white);
    2, the error $\mD = \mC - \Cbar$; 3, its smooth part $\tD$ (red =
    shadow left, blue = overdone); 4, its fine part $\mD - \tD$ (tighter
    scale).  The dot is stain, not error, and saturates on purpose.  Made
    by \code{shading\_review.show\_r3\_maps} (step R3a).
    }
\label{fig:shading-R3a}
\end{figure}

\subsection{Sorted error and sorted stain values}

\begin{figure}[h!]
    \centering
    \includegraphics[width=\textwidth]{shading-R3b}
    \caption{
    C1032 royal\_blue.B: each method's values in one region, sorted from
    lowest to highest, against the percent of pixels below.  A flat curve
    near 0 is a good correction; a steep end is a few pixels far off.  Top:
    $\mD$ on the dye-free inner pad and on the dye-free rim, $\mD - \tD$ on
    the dye-free zone.  Bottom: $\mC$ on the whole pad, the dot, and the
    dye-free zone.  Made by \code{shading\_review.show\_r3\_sorted}
    (step R3b).
    }
\label{fig:shading-R3b}
\end{figure}

\subsection{The numbers}

Standard deviations (sd) in absorbance; ``robust sd'' is $1.4826 \times$ the
median distance from the median.  Made by \code{shading\_review.r3\_numbers}
(step R3c).

\medskip
\begin{center}
\small
\begin{tabular}{@{}lrrrrrrrr@{}}
\toprule
 & \multicolumn{4}{c}{dye-free inner pad (760 px)} &
   \multicolumn{3}{c}{dye-free rim (686 px)} & dot \\
\cmidrule(lr){2-5}\cmidrule(lr){6-8}\cmidrule(l){9-9}
method & sd of $\mD$ & robust & sd of $\tD$ & sd of $\mD - \tD$ & sd of $\mD$
  & sd of $\mD - \tD$ & 99th pct & mean $\mC$ \\
\midrule
$\MA$  & 0.0068 & 0.0066 & 0.0038 & 0.0065 & 0.0411 & 0.0299 & +0.177 & 0.4941 \\
$\MB$1 & \textbf{0.0066} & 0.0061 & 0.0028 & \textbf{0.0058} & 0.0203 & 0.0184 & +0.061 & 0.4934 \\
$\MB$2 & 0.0073 & 0.0071 & 0.0035 & 0.0062 & 0.0296 & 0.0252 & +0.111 & 0.4933 \\
$\MC$  & 0.0074 & 0.0064 & 0.0026 & 0.0067 & 0.0122 & 0.0100 & +0.032 & 0.4935 \\
$\MD$  & 0.0071 & \textbf{0.0060} & \textbf{0.0026} & 0.0063 & \textbf{0.0119} & \textbf{0.0099} & +0.034 & 0.4935 \\
\bottomrule
\end{tabular}
\end{center}

\subsection{What the results show}

\begin{enumerate}
  \item \textbf{No method wins everywhere.}  On the dye-free inner pad,
    $\MB$1 is best (sd 0.0066) and $\MC$ worst (0.0074).  On the dye-free
    rim, $\MC$ and $\MD$ are best (0.012) and $\MB$1 is about twice as bad
    (0.020).
  \item \textbf{The rim matters more.}  Its error is 2--6 times the inner
    pad's under every method.
  \item \textbf{Smoothing the pre-image pulls the edge shadow inward}: $\MB$
    under-corrects the outermost pixel and over-corrects the next ones in.
    At 2 px the effect reaches the inner pad ($\MB$2's $\Cbar$ drops to
    0.0027).
  \item \textbf{$\MD$ is about as good as $\MC$ on the rim and slightly
    better inside}, but its one $\hat\lambda$ is set by the rim (0.85 there,
    0.24 inside).
  \item \textbf{Every method leaves some error at the pad's very edge},
    over-corrected in places and under-corrected in others, because the
    two photos' edges do not line up to the pixel.
  \item On this pad the lined-up pre-image is already smoothed to 0.80
    (rows) and 0.72 (columns) of its scatter (Section~\ref{sec:lineup}).
    On pads that land half-way between pixels, $\MC$ is smoothed much more.
\end{enumerate}

%======================================================================
\section{Side effects: the stain}

The correction must not change the stain.  The dot's mean $\mC$ is
0.4933--0.4941 for all five, a spread of 0.0008 (0.2\,\% of the dot's
value), and its sorted values lie on top of each other
(Figure~\ref{fig:shading-R3b}, bottom middle).  Through the dot, the
methods differ only at the pad's edges (Figure~\ref{fig:shading-R3e},
panel 4).

%======================================================================
\section{How big are these effects?}

Compared with the spreads already known (in absorbance): camera noise
about 0.002, paper texture 0.006--0.008, rim shadow about 0.04.
\begin{itemize}
  \item On the inner pad the methods range from 0.0066 to 0.0074, all
    within the paper's texture.  Choosing among them there is polish.
  \item On the rim they range from 0.012 to 0.041: a large effect.  This is
    where the choice matters.
  \item Lining up lowers the pre-image's scatter by up to half (0.0067 to
    0.0036 on C1020): as large as the differences between methods inside
    the pad.
\end{itemize}

%======================================================================
\section{What this note does not tell us, and what is next}
\label{sec:open}

\begin{itemize}
  \item It is one pad and one channel.  Whether $\hat\lambda$ and the order
    of the methods hold on the other pads is not known; the lining-up
    result says they may not, since each pad's $\mB$ is smoothed by a
    different amount.
  \item \textbf{$\lambda$ (Bill, 2026-10-03).}  A $\lambda$ that varies across
    the pad risks being confused with the stain.  First try to make one
    $\lambda$ work through the right smoothing near the edge.  Only if the
    outermost one to three rows along each pad edge are shaded in a
    fundamentally different way, let $\lambda$ vary there, and only there.
    $\lambda$ may be found once and stored, or set for each photo from its
    noise and how well the two photos line up.
  \item \textbf{The rim} (deferred by Bill on 2026-10-02): why the outermost
    pixel is over-corrected on one side and under-corrected on the other,
    and whether a smoothing that respects the edge can avoid it.  Bill's
    expected method: two-dimensional smoothing on the inner pad and
    smoothing along rows or columns only, at twice the width, near the
    edges, then method $\MD$ with both weights non-zero.
  \item \textbf{Lining up.}  Possible ways to handle the pad-to-pad
    difference in smoothing: record each pad's fractional position and
    allow for it when setting $\sigma$ and $\lambda$; or line up so that the
    smoothing is the same on every pad.  Bill (2026-10-04): compare
    recording the fractional position with lining up by nearest pixel, in
    the rim and edge-smoothing step.
  \item The error is measured against a flat dye-free pad.  Faint stain
    more than 6 px from the dot, if any, would count as error.
\end{itemize}

%======================================================================
\section{How to reproduce}

Script: \file{desktop/padanalysis/shading\_review.py}.  Every figure and
number in this note comes from
\begin{verbatim}
cd desktop/padanalysis
timeout -k 10 600 ../../.venv/bin/python3 -u -c \
  "import shading_review as sr; sr.tex_figs()"
\end{verbatim}
\noindent which writes \file{tex/figs/shading-R2a.pdf}, \file{shading-R3a},
\file{-R3b}, \file{-R3d}, \file{-R3e}, \file{-R4a} and \file{-R4b}, and
prints the numbers.  Functions: \code{post\_pre} (the two photos),
\code{r3\_methods} (the four methods and the error), \code{r3\_numbers},
\code{show\_white}, \code{show\_r3\_maps}, \code{show\_r3\_sorted},
\code{show\_r3\_profile}, \code{r3\_example} (worked example; prints the
code's values relative to $\bw$), \code{lineup\_noise},
\code{lineup\_noise\_table}, \code{show\_lineup\_noise},
\code{show\_lineup\_ratio}.  In the code, the arrays \code{A\_post} and
\code{A\_pre} are $\mA - \bw$ and $\mB - \bw$, \code{corr} is $\hB - \bw$,
\code{after} is $\mC$, \code{level} is $\Cbar$, \code{err} is $\mD$, and
\code{wt} is $\hat\lambda$.  The study log is
\file{docs/shading\_study\_2026-09-29.md}, section ``Reassessment''.

\end{document}
